% =========================================================================
% Capítulo 0: Repaso
% Tema 02: Cálculo vectorial
% Problema 03: Derivadas de Vectores y Productos
% =========================================================================
\subsubsection{Derivadas de Vectores y Productos}
Considere que $\vec{u}(t)$ y $\vec{v}(t)$ son vectores que dependen del tiempo. Calcule:

\begin{enumerate}
    \item Halle $\frac{d}{dt} (\vec{u} \cdot \vec{v})$.
    \item Halle $\frac{d}{dt} (\vec{u} \times \vec{v})$.
    \item Sea $\vec{r}(t) = \cos(\omega t)\vec{i} + \sin(\omega t)\vec{j} + bt\vec{k}$. Calcule $\dot{\vec{r}}$ y $\ddot{\vec{r}}$.
    \item Si $\vec{v} = \dot{\vec{r}}$, calcule la derivada de la energía cinética escalar: $\frac{d}{dt} (\frac{1}{2} m \vec{v} \cdot \vec{v})$.
\end{enumerate}
% --- SOLUCIÓN ---------------------------------------------------------------
\begin{solucion}{Derivación Vectorial}
\begin{enumerate}
    \item $\frac{d}{dt} (\vec{u} \cdot \vec{v}) = \dot{\vec{u}} \cdot \vec{v} + \vec{u} \cdot \dot{\vec{v}}$.
    \item $\vec{r}(t) = \cos(\omega t)\vec{i} + \sin(\omega t)\vec{j} + bt\vec{k}$. 
    Derivando componente a componente:
    $\dot{\vec{r}} = -\omega\sin(\omega t)\vec{i} + \omega\cos(\omega t)\vec{j} + b\vec{k}$.
    $\ddot{\vec{r}} = -\omega^2\cos(\omega t)\vec{i} - \omega^2\sin(\omega t)\vec{j}$.
    \item $\frac{d}{dt} (\frac{1}{2} m \vec{v} \cdot \vec{v}) = \frac{1}{2} m (\dot{\vec{v}} \cdot \vec{v} + \vec{v} \cdot \dot{\vec{v}}) = m \vec{v} \cdot \dot{\vec{v}} = m \vec{v} \cdot \vec{a} = P$ (Potencia).
\end{enumerate}
\end{solucion}
